\documentclass[a4paper,12pt]{article}

% Packages
\usepackage[utf8]{inputenc}
\usepackage[T1]{fontenc}
\usepackage[french]{babel}
\usepackage{amsmath, amssymb, amsthm}
\usepackage{geometry}
\usepackage{xcolor}
\usepackage{tcolorbox}
\usepackage{enumitem}
\usepackage{fancyhdr}
\usepackage{array}
\usepackage{booktabs}
\usepackage{tikz}
\usepackage{circuitikz}
\usepackage{multirow}
\usepackage{graphicx}

\usetikzlibrary{shapes.gates.logic.US,positioning}

% Configuration de la page
\geometry{margin=2.5cm}
\pagestyle{fancy}
\fancyhf{}
\fancyhead[L]{Fondements de l'Informatique \& Structure des Ordinateurs}
\fancyhead[R]{Licence Informatique - S1 \& S2}
\fancyfoot[C]{\thepage}

% Environnements personnalisés
\tcbuselibrary{theorems}

\newtcolorbox{defbox}[1]{
    colback=blue!5!white,
    colframe=blue!75!black,
    fonttitle=\bfseries,
    title=#1
}

\newtcolorbox{propbox}[1]{
    colback=green!5!white,
    colframe=green!75!black,
    fonttitle=\bfseries,
    title=#1
}

\newtcolorbox{theorembox}[1]{
    colback=red!5!white,
    colframe=red!75!black,
    fonttitle=\bfseries,
    title=#1
}

\newtcolorbox{examplebox}[1]{
    colback=orange!5!white,
    colframe=orange!75!black,
    fonttitle=\bfseries,
    title=#1
}

\newtcolorbox{remarkbox}[1]{
    colback=purple!5!white,
    colframe=purple!75!black,
    fonttitle=\bfseries,
    title=#1
}

% Titre
\title{\textbf{Fondements de l'Informatique \\ et \\ Structure des Ordinateurs} \\ 
\large L'Algèbre de Boole, la Représentation de l'Information \\ et l'Architecture des Systèmes Numériques}
\author{Arthur LOUVET}
\date{Licence Informatique - Semestres 1 \& 2 (2024-2025)}

\begin{document}

\maketitle
\thispagestyle{empty}

\newpage
\tableofcontents
\newpage

% ============================================================================
% PARTIE I : ALGÈBRE DE BOOLE ET FONCTIONS LOGIQUES
% ============================================================================

\part{Algèbre de Boole et Fonctions Logiques}

\section{L'Algèbre de Boole}

\subsection{Le domaine de définition}

\begin{defbox}{Définition : Ensemble de définition des nombres booléens}
L'ensemble de définition des nombres booléens est $\{0,1\}$ mais ces nombres ne peuvent prendre que l'existence : 
\begin{itemize}
    \item $0 \rightarrow$ valeur fausse
    \item $1 \rightarrow$ valeur vraie
\end{itemize}
\end{defbox}

\subsection{Les opérateurs de base}

\subsubsection{Opérations unaires et binaires}

\begin{defbox}{Définition : Opérations de base}
Ces nombres utilisent 4 opérateurs :
\begin{itemize}
    \item \textbf{NON} (non $a$, $-a$, $\overline{a}$) $\rightarrow$ opération unaire \\
    $a \neq b \Rightarrow 1 = 0$
    
    \item \textbf{OU} ($a \vee b$, $a + b$) $\rightarrow$ opération binaire
    
    \item \textbf{ET} ($a \wedge b$, $a \cdot b$) $\rightarrow$ opération binaire
\end{itemize}
\end{defbox}

\subsection{Propriétés fondamentales}

\begin{propbox}{Propriétés fondamentales de l'algèbre de Boole}
Ces nombres possèdent des propriétés fondamentales :

\begin{enumerate}[leftmargin=2cm]
    \item \textbf{Idempotence :} $a + a = a$ ~~/~~ $a \cdot a = a$
    
    \item \textbf{Involution :} $\overline{\overline{a}} = a$
    
    \item \textbf{Commutativité :} $a + b = b + a$ ~~/~~ $a \cdot b = b \cdot a$
    
    \item \textbf{Élément neutre :} $a + 0 = a$ ~~/~~ $a \cdot 1 = a$
    
    \item \textbf{Élément absorbant :} $a + 1 = 1$ ~~/~~ $a \cdot 0 = 0$
    
    \item \textbf{Complémentarité :} $a + \overline{a} = 1$ ~~/~~ $a \cdot \overline{a} = 0$
    
    \item \textbf{Associativité :} $(a + b) + c = a + (b + c)$ \\
    $(a \cdot b) \cdot c = a \cdot (b \cdot c)$
\end{enumerate}
\end{propbox}

\begin{propbox}{Propriétés de distributivité}
\begin{itemize}
    \item \textbf{Distributivité de ``$\vee$'' sur ``$\wedge$'' :}
    \[(a + b) \cdot c = (a \cdot c) + (b \cdot c)\]
    
    \item \textbf{Distributivité de ``$\wedge$'' sur ``$\vee$'' :}
    \[(a \cdot b) + c = (a + c) \cdot (b + c)\]
\end{itemize}
\end{propbox}

\subsection{Règle de De Morgan}

\begin{theorembox}{Théorème : Lois de De Morgan}
Pour tous $a, b \in \{0,1\}$ :

\begin{align*}
\overline{(a + b)} &= \overline{a} \cdot \overline{b} \\[0.5em]
\overline{(a \cdot b)} &= \overline{a} + \overline{b}
\end{align*}

Ces lois permettent de transformer la négation d'une somme en produit de négations, et réciproquement.
\end{theorembox}

\newpage
\section{Les autres opérateurs usuels}

\subsection{Le OU booléen (XOR)}

\begin{defbox}{Définition : OU exclusif (XOR)}
Le OU booléen (XOR) noté $a \oplus b$ est différent du OU classique car il est seulement vrai lorsque les deux valeurs sont différentes. Pour que $a \oplus b$ soit vraie, il faut exactement l'une des deux valeurs soit vraie (mais pas les deux).
\end{defbox}

\vspace{0.5cm}

\begin{center}
\textbf{Table de vérité du OU exclusif}
\vspace{0.3cm}

\begin{tabular}{|c|c|c|c|}
\hline
$a$ & $b$ & $a \vee b$ & $a \oplus b$ \\
\hline
0 & 0 & 0 & 0 \\
\hline
0 & 1 & 1 & 1 \\
\hline
1 & 0 & 1 & 1 \\
\hline
1 & 1 & 1 & 0 \\
\hline
\end{tabular}
\end{center}

\subsection{Le NON OU (NOR)}

\begin{defbox}{Définition : NOR}
Le ``non OU'' (NOR) noté $\text{NOR}(a,b)$ est différent du OU classique car il est seulement vrai lorsque $a$ et $b$ sont faux (que celui-ci soit vrai).
\end{defbox}

\vspace{0.5cm}

\begin{center}
\textbf{Table de vérité du NOR}
\vspace{0.3cm}

\begin{tabular}{|c|c|c|c|}
\hline
$a$ & $b$ & $a \vee b$ & $\text{NOR}$ \\
\hline
0 & 0 & 0 & 1 \\
\hline
0 & 1 & 1 & 0 \\
\hline
1 & 0 & 1 & 0 \\
\hline
1 & 1 & 1 & 0 \\
\hline
\end{tabular}
\end{center}

\subsection{Le NON ET (NAND)}

\begin{defbox}{Définition : NAND}
Le ``non ET'' (NAND) se note $\overline{a \cdot b}$ et est différent du ET car pour le ET, les valeurs doivent être vraies. Alors que pour le non ET, les valeurs ne doivent pas être vraies en même temps.
\end{defbox}

\vspace{0.5cm}

\begin{center}
\textbf{Table de vérité du NAND}
\vspace{0.3cm}

\begin{tabular}{|c|c|c|c|}
\hline
$a$ & $b$ & $a \cdot b$ & $\overline{a \cdot b}$ \\
\hline
0 & 0 & 0 & 1 \\
\hline
0 & 1 & 0 & 1 \\
\hline
1 & 0 & 0 & 1 \\
\hline
1 & 1 & 1 & 0 \\
\hline
\end{tabular}
\end{center}

\newpage
\section{Fonctions logiques}

\subsection{Propriétés des opérateurs}

\subsubsection{La commutativité}

\begin{propbox}{Propriété : Rappel pour tous $a,b$}
\begin{itemize}
    \item $a + b = b + a$
    \item $a \cdot b = b \cdot a$
\end{itemize}
\end{propbox}

\subsubsection{L'associativité}

\begin{propbox}{Propriété : Attention pour NOR et NAND}
Attention : $\forall a,b,c \in \{0,1\}$ :
\begin{itemize}
    \item $\overline{a + b + c} \neq (\overline{a + b}) + c$
    \item $\overline{a \cdot b \cdot c} \neq (\overline{a \cdot b}) \cdot c$
\end{itemize}

Ainsi, il n'y a aucune propriété / manière évidente en NOR et NAND (il n'y a pas de propriété d'associativité).
\end{propbox}

\begin{examplebox}{Remarque}
On associe couramment le ``NOR'' à d'autres calculs mais le principe est le même !!!
\end{examplebox}

\subsection{Groupes et opérateur complet}

\subsubsection{L'implication}

\begin{defbox}{Définition : Implication}
Voici la table de vérité de l'implication :

\vspace{0.3cm}
\begin{center}
\begin{tabular}{|c|c|c|c|}
\hline
$a$ & $b$ & $a \Rightarrow b$ & $b \Rightarrow a$ \\
\hline
0 & 0 & 1 & 1 \\
\hline
0 & 1 & 1 & 0 \\
\hline
1 & 0 & 0 & 1 \\
\hline
1 & 1 & 1 & 1 \\
\hline
\end{tabular}
\end{center}
\end{defbox}

\subsubsection{Opérateur / Groupe et opérateur complet}

\begin{defbox}{Définition : Groupe complet}
L'ensemble $\{+, \cdot, -\}$ est un groupe complet.

L'ensemble $\{\oplus, -\}$ est-il complet ?

\vspace{0.3cm}
Pour tous $a, b \in \{0,1\}$ :
\begin{itemize}
    \item $a \cdot b = \overline{a} \oplus b$ (pour le ET)
    \item $a + b = \overline{(\overline{a + b})} = \overline{(\overline{a} \cdot \overline{b})}$
\end{itemize}

Donc $\{\oplus, -\}$ est donc bel et bien un ensemble complet.
\end{defbox}

\begin{propbox}{Propriété}
Les fonctions logiques ou booléennes (+ pour canonique) :

Soit $f_1(a,b,c) = a + b$ et $f_2(a,b,c) = (a \cdot c) + (\overline{a} \cdot b) + a \cdot c$

$f_2(a,b,c) = \overline{a} \cdot b \cdot c + a \cdot \overline{b} + \overline{a} \cdot \overline{b} \cdot \overline{c}$
\end{propbox}

\newpage
\section{Les égalités de fonction}

\subsection{Vérification par table de vérité}

\begin{examplebox}{Exemple avec table de vérité}
Voici une table de vérité complète pour fonctions à 3 variables :

\vspace{0.5cm}
\begin{center}
\begin{tabular}{|c|c|c|c|c|c|c|c|c|c|}
\hline
$a$ & $b$ & $c$ & $\overline{a}$ & $\overline{b}$ & $a \cdot \overline{b}$ & $\overline{a} \cdot \overline{b}$ & $a \cdot b$ & $c + a$ & $f_c$ \\
\hline
0 & 0 & 0 & 1 & 1 & 0 & 1 & 0 & 0 & 0 \\
\hline
0 & 0 & 1 & 1 & 1 & 0 & 1 & 0 & 1 & 0 \\
\hline
0 & 1 & 0 & 1 & 0 & 0 & 0 & 0 & 0 & 0 \\
\hline
0 & 1 & 1 & 1 & 0 & 0 & 0 & 0 & 1 & 0 \\
\hline
1 & 0 & 0 & 0 & 1 & 1 & 0 & 0 & 1 & 1 \\
\hline
1 & 0 & 1 & 0 & 1 & 1 & 0 & 0 & 1 & 0 \\
\hline
1 & 1 & 0 & 0 & 0 & 0 & 0 & 1 & 1 & 1 \\
\hline
1 & 1 & 1 & 0 & 0 & 0 & 0 & 1 & 1 & 1 \\
\hline
\end{tabular}
\end{center}
\end{examplebox}

\subsection{Normalisation de l'écriture d'une fonction}

\begin{defbox}{Définition : Formule normale}
La formule normale se compose de deux formes :

\textbf{Forme disjonctive :}
\begin{itemize}
    \item Les opérateurs conjonctifs ne portent que sur des variables
    \item Les variables sont regroupées par forme : conjonction
    \item Les termes sont regroupés : $\left[\frac{\text{OU}}{\text{ET}}\right]$ ou $\left[\frac{\text{OU}}{?}\right]$
\end{itemize}

\textbf{Forme conjonctive :}
\begin{itemize}
    \item Les opérateurs disjonctifs ne portent que sur des variables
    \item Les variables sont regroupées par formes : disjonction
    \item Les termes sont regroupés : $\left[\frac{\text{ET}}{\text{OU}}\right]$ ou $\left[\frac{\text{ET}}{?}\right]$
\end{itemize}

\vspace{0.3cm}
Exemples :
\begin{itemize}
    \item \textbf{FND :} $f(a,b,c) = (a \cdot \overline{b} \cdot c) + (\overline{a} \cdot b)$
    \item \textbf{FNC :} $f(a,b,c) = (a + b + \overline{c}) \cdot (b + c)$
\end{itemize}
\end{defbox}

\subsection{La forme canonique : Passage entre la formule}

\begin{defbox}{Définition : Forme canonique}
\textbf{Forme canonique disjonctive :}
\[f_1(a,b,c) = \overline{a}\overline{b}c + \overline{a}b\overline{c} + abc \quad \text{(disjonctive)}\]

\textbf{Forme canonique conjonctive :}
\[f_2(a,b,c) = (a+b+\overline{c}) \cdot (a+\overline{b}+c) \cdot (\overline{a}+\overline{b}+c) \quad \text{(conjonctive)}\]
\end{defbox}

\begin{examplebox}{Exemple : Table de vérité}
\begin{center}
\begin{tabular}{|c|c|c|c|c|}
\hline
$a$ & $b$ & $c$ & $f_1$ & $f_2$ \\
\hline
0 & 0 & 0 & 0 & 1 \\
\hline
0 & 0 & 1 & 1 & 0 \\
\hline
0 & 1 & 0 & 1 & 0 \\
\hline
0 & 1 & 1 & 0 & 1 \\
\hline
1 & 0 & 0 & 0 & 1 \\
\hline
1 & 0 & 1 & 0 & 0 \\
\hline
1 & 1 & 0 & 0 & 1 \\
\hline
1 & 1 & 1 & 1 & 0 \\
\hline
\end{tabular}
\end{center}

\vspace{0.3cm}
Comment passer à une forme canonique ?

\[f_c = (a+b+c) \cdot (a+\overline{c}) \cdot (\overline{a}+\overline{b}+c)\]
\end{examplebox}

\newpage
\section{Fonctions incomplètes et états indéterminés}

\subsection{Principe}

\begin{defbox}{Définition : États indéterminés}
Pour la transformation d'états indéterminés, on s'en sert pour simplifier les fonctions en choisissant le $1$ le plus souvent possible.

Lorsque l'on a plus de bits que nécessaires, ceux en trop sont dans un état indéterminé. C'est une \textbf{optimisation de bits et de place}.
\end{defbox}

\begin{remarkbox}{Remarque importante}
Voici pourquoi on simplifie les fonctions : pour les récupérer et les « économiser ».
\end{remarkbox}

\subsection{Exemple de fonctions incomplètes}

Soient les fonctions suivantes :

\begin{align}
M_1 &= V_1 V_2 + V_3 V_4 \\
M_2 &= V_2 V_4 + V_1 V_3 \\
M_3 &= V_3 V_4 + V_2 V_3
\end{align}

\subsection{Table de vérité avec états indéterminés}

\begin{table}[h]
\centering
\begin{tabular}{cccc|ccc}
\toprule
$V_1$ & $V_2$ & $V_3$ & $V_4$ & $M_1$ & $M_2$ & $M_3$ \\
\midrule
0 & 0 & 0 & 0 & 0 & 0 & 0 \\
0 & 0 & 0 & 1 & - & 0 & 0 \\
0 & 0 & 1 & 0 & 0 & 0 & 1 \\
0 & 0 & 1 & 1 & 1 & 0 & 1 \\
0 & 1 & 0 & 0 & 0 & 0 & 0 \\
0 & 1 & 0 & 1 & 0 & 1 & 0 \\
0 & 1 & 1 & 0 & 0 & 0 & 1 \\
0 & 1 & 1 & 1 & 1 & 1 & 1 \\
1 & 0 & 0 & 0 & 0 & 0 & 0 \\
1 & 0 & 0 & 1 & - & 1 & 0 \\
1 & 0 & 1 & 0 & 0 & 1 & 1 \\
1 & 0 & 1 & 1 & 1 & 1 & 1 \\
1 & 1 & 0 & 0 & 1 & 0 & 0 \\
1 & 1 & 0 & 1 & 1 & 1 & 0 \\
1 & 1 & 1 & 0 & 1 & 1 & 1 \\
1 & 1 & 1 & 1 & 1 & 1 & 1 \\
\bottomrule
\end{tabular}
\caption{Table de vérité avec états indéterminés (notés « - »)}
\end{table}

\subsection{Formes normales}

\begin{itemize}
    \item \textbf{FND (Forme Normale Disjonctive)} : SOMME = 1
    \item \textbf{FNC (Forme Normale Conjonctive)} : SOMME = 0
\end{itemize}

\section{La simplification d'une fonction}

Soit :
\begin{align*}
f_1 &= abc + b(a+\overline{c}) = abc + ab + b\overline{c} \\
f_1 &= abc + abc + ab\overline{c} + ab\overline{c} + a\overline{b}\overline{c} + \overline{a}b\overline{c} \\
f_1 &= abc + ab\overline{c} + abc + a\overline{b}\overline{c} = ab + b\overline{c}
\end{align*}

\newpage

% ============================================================================
% PARTIE II : REPRÉSENTATION DE L'INFORMATION
% ============================================================================

\part{Représentation de l'Information}

\section{Les entiers naturels (taille infinie) $\mathbb{N}$}

\subsection{Types de nombres}

\begin{itemize}
    \item Relatifs $\infty$ : $\mathbb{Z}$ (complément à la base)
    \item Décimaux $\infty$ : $\mathbb{D}$
    \item Rationnels : arithmétique infinie $\mathbb{Q}$
    \item Réels : $\mathbb{R}$
    \item Complexe : $\mathbb{C}$
\end{itemize}

\subsection{Le codage}

\begin{defbox}{Définition : Codage N}
Le codage $N = \sum_{i=0}^{k} n_i \cdot b^i$ où :
\begin{itemize}
    \item $k$ : nombre de poids de chiffres à gauche
    \item chiffres de gauche, rangés en groupe $k^{\text{ième}}$
    \item $b$ : base de numération, $\geq 2$ arbitraire (complément à la base)
\end{itemize}
\end{defbox}

\begin{examplebox}{Exemple : Langage européen (base 10)}
\[566 = 5 \times 10^2 + 6 \times 10^1 + 6 \times 10^0\]

\vspace{0.3cm}
En base 9 :
\begin{align*}
(568)_9 &= (5 \times 9^2 + 6 \times 9^1 + 8 \times 9^0) = 405 + 54 + 8 = (467)_{10} \\
(5685)_9 &= (5 \times 9^3 + 6 \times 9^2 + 8 \times 9^1 + 5 \times 9^0) = (4172)_{10}
\end{align*}
\end{examplebox}

\subsection{Conversion : division euclidienne}

\begin{examplebox}{Exemple : Conversion de base}
\begin{align*}
\text{Exemple : } (10111101)_2 &= (189)_{10}
\end{align*}

\vspace{0.3cm}
Si on veut convertir 585 en base 7, on utilise la division euclidienne :

\[
\begin{array}{rcl}
585 &=& 7 \times 83 + 4 \\
83 &=& 7 \times 11 + 6 \\
11 &=& 7 \times 1 + 4 \\
1 &=& 7 \times 0 + 1
\end{array}
\]

Résultats : $585_{10} = (1464)_7$
\end{examplebox}

\subsubsection{Principe général}

\begin{align*}
(N)_b &= \sum n_i \cdot b^i \\
N &= n_0 + b(n_1 + b(n_2 + b \cdot n_3 + \ldots)) \\
N &= b \times Q + R
\end{align*}

où $R$ est le reste de la division (chiffre de poids faible) et $Q$ est le quotient.

\newpage
\section{Systèmes de codage}

\subsection{Rappel : système binaire}

Le binaire est un \textbf{code pondéré}. La formule binomiale permet de convertir un nombre binaire en décimal :

\[
\text{Valeur décimale} = \sum_{i=0}^{n-1} B_i \times 2^i
\]

où $B_i$ est le bit de rang $i$.

\subsection{Codes pondérés}

\subsubsection{Code BCD (Binary Coded Decimal)}

\begin{defbox}{Définition : Code BCD}
\begin{itemize}
    \item Code uniquement les 10 chiffres décimaux (0 à 9)
    \item Chaque chiffre décimal est représenté par 4 bits
    \item Code pondéré et réfléchi
\end{itemize}
\end{defbox}

\subsubsection{Code de Aiken}

\begin{defbox}{Définition : Code de Aiken}
Le code de Aiken possède les propriétés suivantes :

\begin{itemize}
    \item Pratique
    \item Pondéré avec les poids : \textbf{2, 4, 2, 1}
    \item BCDique (code les 10 chiffres)
    \item Réfléchi
    \item Autocomplémentaire
\end{itemize}
\end{defbox}

\begin{table}[h]
\centering
\begin{tabular}{c|cccc|c}
\toprule
\textbf{Décimal} & \multicolumn{4}{c|}{\textbf{Poids}} & \textbf{Code de Aiken} \\
 & 2 & 4 & 2 & 1 & \\
\midrule
0 & 0 & 0 & 0 & 0 & 0000 \\
1 & 0 & 0 & 0 & 1 & 0001 \\
2 & 0 & 0 & 1 & 0 & 0010 \\
3 & 0 & 0 & 1 & 1 & 0011 \\
4 & 0 & 1 & 0 & 0 & 0100 \\
5 & 1 & 0 & 1 & 1 & 1011 \\
6 & 1 & 1 & 0 & 0 & 1100 \\
7 & 1 & 1 & 0 & 1 & 1101 \\
8 & 1 & 1 & 1 & 0 & 1110 \\
9 & 1 & 1 & 1 & 1 & 1111 \\
\bottomrule
\end{tabular}
\caption{Table du code de Aiken}
\end{table}

\subsection{Codes non-pondérés}

\subsubsection{Code excédant 3 (+3)}

\begin{defbox}{Définition : Code excédant 3}
\begin{itemize}
    \item Code BCDique
    \item On prend le code binaire naturel et on ajoute 3
    \item Exemple : $0 \rightarrow 0+3 = 3$ (0011), $1 \rightarrow 1+3 = 4$ (0100), etc.
\end{itemize}
\end{defbox}

\begin{table}[h]
\centering
\begin{tabular}{c|c|c}
\toprule
\textbf{Décimal} & \textbf{Binaire naturel} & \textbf{Excédant 3} \\
\midrule
0 & 0000 & 0011 \\
1 & 0001 & 0100 \\
2 & 0010 & 0101 \\
3 & 0011 & 0110 \\
4 & 0100 & 0111 \\
5 & 0101 & 1000 \\
6 & 0110 & 1001 \\
7 & 0111 & 1010 \\
8 & 1000 & 1011 \\
9 & 1001 & 1100 \\
\bottomrule
\end{tabular}
\caption{Table du code excédant 3}
\end{table}

\newpage
\subsubsection{Code de GRAY}

\begin{defbox}{Définition : Code de GRAY}
Le code de GRAY est un code cyclique très utilisé dans les encodeurs rotatifs et les dispositifs électromécaniques.

\textbf{Principe de construction :}

\begin{enumerate}
    \item On commence à 0
    \item Pour la valeur suivante, on inverse le bit le plus à droite
    \item On ne retombe jamais sur une valeur déjà connue
    \item Un seul bit change entre deux valeurs consécutives
\end{enumerate}

\textbf{Avantages :}

\begin{itemize}
    \item Évite les erreurs de lecture lors de transitions
    \item Particulièrement utile pour les roues/molettes électromécaniques
    \item Un seul bit change à la fois = moins d'erreurs de lecture
\end{itemize}
\end{defbox}

\begin{table}[h]
\centering
\begin{tabular}{c|c|c}
\toprule
\textbf{Décimal} & \textbf{Binaire naturel} & \textbf{Code de GRAY} \\
\midrule
0 & 0000 & 0000 \\
1 & 0001 & 0001 \\
2 & 0010 & 0011 \\
3 & 0011 & 0010 \\
4 & 0100 & 0110 \\
5 & 0101 & 0111 \\
6 & 0110 & 0101 \\
7 & 0111 & 0100 \\
8 & 1000 & 1100 \\
9 & 1001 & 1101 \\
10 & 1010 & 1111 \\
11 & 1011 & 1110 \\
12 & 1100 & 1010 \\
13 & 1101 & 1011 \\
14 & 1110 & 1001 \\
15 & 1111 & 1000 \\
\bottomrule
\end{tabular}
\caption{Table du code de GRAY (4 bits)}
\end{table}

\begin{remarkbox}{Application pratique}
Le code de GRAY est utilisé dans les encodeurs rotatifs (comme les molettes de souris, les potentiomètres numériques, etc.) car lors de la rotation, un seul segment change d'état à la fois, ce qui évite les erreurs de lecture dues aux désynchronisations mécaniques.
\end{remarkbox}

\newpage
\section{Codes redondants - Détection d'erreurs}

\subsection{Bit de parité}

\begin{defbox}{Définition : Bit de parité}
Le bit de parité permet d'effectuer des tests de communication entre machines et de détecter des erreurs de transmission.

Pour déterminer la parité d'une valeur en binaire, on fait un \textbf{OU EXCLUSIF (XOR)} entre chaque valeur de manière cumulative (au fur et à mesure).
\end{defbox}

\begin{examplebox}{Exemple : Calcul de parité}
Soit la valeur binaire : $1011$

\begin{align*}
\text{Parité} &= 1 \oplus 0 \oplus 1 \oplus 1 \\
              &= 1 \oplus 0 = 1 \\
              &= 1 \oplus 1 = 0 \\
              &= 0 \oplus 1 = 1
\end{align*}

Le bit de parité est donc $1$ (parité impaire).
\end{examplebox}

\begin{remarkbox}{Application}
Exemple : barrette mémoire avec détection d'erreurs utilisant ce fonctionnement.
\end{remarkbox}

\subsection{Code pondéré redondant}

\begin{defbox}{Définition : Code redondant à 5 bits}
Pour le code pondéré redondant, on utilise 5 bits ($2^5 = 32$ combinaisons possibles).

Propriété importante : chaque code de valeurs ne contient que des 1, donc s'il n'y a pas 2 × 1, on peut détecter une erreur.
\end{defbox}

\subsection{Code ASCII}

\begin{defbox}{Définition : Code ASCII}
Le code \textbf{ASCII} (American Standard Code for Information Interchange) est :
\begin{itemize}
    \item Codé sur $2^7$ bits (128 valeurs)
    \item + 1 bit de parité (8 bits au total)
    \item ou ce bit utilisé pour un code ASCII étendu (variable selon les pays pour les caractères différents)
\end{itemize}
\end{defbox}

\newpage

% ============================================================================
% PARTIE III : CIRCUITS ET ARCHITECTURE
% ============================================================================

\part{Circuits et Architecture des Ordinateurs}

\section{Introduction à l'architecture des ordinateurs}

\subsection{Composants de base}

\begin{defbox}{Définition : Composants fondamentaux}
Un ordinateur est composé des éléments fondamentaux suivants :

\begin{itemize}
    \item \textbf{UAL (Unité Arithmétique et Logique)} : ALU en anglais, effectue les opérations arithmétiques et logiques
    \item \textbf{Unité de contrôle} : coordonne les opérations de l'ordinateur
    \item \textbf{Entrées / Sorties} : permettent la communication avec l'extérieur
\end{itemize}
\end{defbox}

\subsection{Opérations binaires}

Les opérations de base se font sur des opérandes binaires $\{0, 1\}$ avec les opérateurs : $+$, $-$, $\times$.

Ces opérations possèdent des propriétés et peuvent être représentées par des fonctions et des tables de vérité, qui permettent ensuite la simplification des circuits logiques.

\section{Les circuits : Transcodeurs}

\subsection{Définition générale}

\begin{defbox}{Définition : Transcodeur}
Dans un circuit, on a le \textbf{transcodeur} : circuit dans lequel on a :
\begin{itemize}
    \item le \textbf{codeur} (pas de code $\rightarrow$ code)
    \item le \textbf{décodeur} (code $\rightarrow$ pas de code)
\end{itemize}

Le transcodeur fait donc : code $\rightarrow$ code
\end{defbox}

\subsection{Le codeur}

\begin{defbox}{Définition : Codeur}
\textbf{Côté CODEUR :} 1 et 1 seule entrée à 1

\begin{itemize}
    \item ``pas de code'' : 4 valeurs, on a 4 touches : 0, 1, 2, 3
    \item Si on appuie sur la première touche, le circuit de codage doit générer 0, etc.
    \item On fait un circuit avec 4 entrées, 2 sorties
    \item Tableau avec $2^4 = 16$ lignes, 2 fonctions
    \item Toutes les autres valeurs que 0, 1, 2 et 3 ne nous intéressent pas (états indéterminés)
\end{itemize}

\textbf{En sortie :} le codeur sort la valeur qui nous intéresse, celle que l'on a mis à 1 (pression).
\end{defbox}

\begin{examplebox}{Exemple : Codeur 4 vers 2}
\begin{center}
\begin{tabular}{|c|cccc|cc|}
\hline
\textbf{Valeur} & $E_3$ & $E_2$ & $E_1$ & $E_0$ & $S_1$ & $S_0$ \\
\hline
0 & 0 & 0 & 0 & 1 & 0 & 0 \\
1 & 0 & 0 & 1 & 0 & 0 & 1 \\
2 & 0 & 1 & 0 & 0 & 1 & 0 \\
3 & 1 & 0 & 0 & 0 & 1 & 1 \\
\hline
\end{tabular}
\end{center}

Les équations logiques sont :
\begin{align*}
S_1 &= E_2 + E_3 \\
S_0 &= E_1 + E_3
\end{align*}
\end{examplebox}

\subsection{Le décodeur}

\begin{defbox}{Définition : Décodeur}
\textbf{Côté DÉCODEUR :}
\begin{itemize}
    \item Prend en entrée 1 valeur codée dans un code
    \item Prend en sortie, la sortie activée dont le numéro correspond à la valeur d'entrée
\end{itemize}

Exemple de décodeur binaire sur 3 bits, 8 sorties.
\end{defbox}

\begin{examplebox}{Exemple : Décodeur 3 vers 8}
\begin{center}
\begin{tabular}{|ccc|cccccccc|}
\hline
$E_2$ & $E_1$ & $E_0$ & $S_0$ & $S_1$ & $S_2$ & $S_3$ & $S_4$ & $S_5$ & $S_6$ & $S_7$ \\
\hline
0 & 0 & 0 & 1 & 0 & 0 & 0 & 0 & 0 & 0 & 0 \\
0 & 0 & 1 & 0 & 1 & 0 & 0 & 0 & 0 & 0 & 0 \\
0 & 1 & 0 & 0 & 0 & 1 & 0 & 0 & 0 & 0 & 0 \\
0 & 1 & 1 & 0 & 0 & 0 & 1 & 0 & 0 & 0 & 0 \\
1 & 0 & 0 & 0 & 0 & 0 & 0 & 1 & 0 & 0 & 0 \\
1 & 0 & 1 & 0 & 0 & 0 & 0 & 0 & 1 & 0 & 0 \\
1 & 1 & 0 & 0 & 0 & 0 & 0 & 0 & 0 & 1 & 0 \\
1 & 1 & 1 & 0 & 0 & 0 & 0 & 0 & 0 & 0 & 1 \\
\hline
\end{tabular}
\end{center}

Les équations logiques sont :
\begin{align*}
S_0 &= \overline{E_2} \cdot \overline{E_1} \cdot \overline{E_0} \\
S_1 &= \overline{E_2} \cdot \overline{E_1} \cdot E_0 \\
S_2 &= \overline{E_2} \cdot E_1 \cdot \overline{E_0} \\
&\text{etc...}
\end{align*}
\end{examplebox}

\subsection{Le transcodeur}

\begin{defbox}{Définition : Transcodeur}
\textbf{Côté TRANSCODEUR :}
\begin{itemize}
    \item Prend en entrée 1 valeur dans 1 code
    \item Prend en sortie la même valeur dans un autre code
\end{itemize}

Exemple : code GRAY sur 3 bits $\rightarrow$ binaire
\end{defbox}

\begin{examplebox}{Exemple : Transcodeur GRAY vers Binaire}
\begin{center}
\begin{tabular}{|c|ccc|ccc|}
\hline
\textbf{Val} & $G_2$ & $G_1$ & $G_0$ & $B_2$ & $B_1$ & $B_0$ \\
\hline
0 & 0 & 0 & 0 & 0 & 0 & 0 \\
1 & 0 & 0 & 1 & 0 & 0 & 1 \\
2 & 0 & 1 & 1 & 0 & 1 & 0 \\
3 & 0 & 1 & 0 & 0 & 1 & 1 \\
4 & 1 & 1 & 0 & 1 & 0 & 0 \\
5 & 1 & 1 & 1 & 1 & 0 & 1 \\
6 & 1 & 0 & 1 & 1 & 1 & 0 \\
7 & 1 & 0 & 0 & 1 & 1 & 1 \\
\hline
\end{tabular}
\end{center}

Les équations de conversion sont :
\begin{align*}
B_2 &= G_2 \\
B_1 &= G_2 \oplus G_1 \\
B_0 &= G_2 \oplus G_1 \oplus G_0
\end{align*}
\end{examplebox}

\newpage
\section{Circulation de l'information - Sur les routes de l'information}

\subsection{Introduction}

\begin{defbox}{Définition : Communication dans un processeur}
Un processeur est composé de plusieurs circuits remplissant chacun leurs fonctions. Entre ces circuits, les informations \textbf{circulent} sur des \textbf{bus}. 

Le terme \textbf{bus} désigne l'ensemble du moyen de transport, matériel et logiciel (protocole).
\end{defbox}

\subsection{Bus série et bus parallèle}

Le terme information désigne ici une valeur quelconque, codée sur un certain nombre de bits. La question de la simultanéité de la circulation des bits se pose.

\subsubsection{Bus série}

\begin{defbox}{Définition : Bus série}
On parle de communication et de \textbf{bus série} quand les bits circulent \textbf{les uns après les autres} sur un même support.

\vspace{0.3cm}
\textbf{Exemples :} RS232, USB, SATA, Ethernet
\end{defbox}

\begin{center}
\textbf{Représentation d'un bus série}

\vspace{0.5cm}
\begin{tikzpicture}
\draw[thick, ->] (0,0) -- (10,0);
\foreach \x/\bit in {0/1, 1/0, 2/1, 3/1, 4/0, 5/0, 6/1, 7/0} {
    \node at (\x+0.5, 0.4) {\small \bit};
}
\node[right] at (10.2,0) {temps $\rightarrow$};
\node at (5,-1) {... 10110010 ...};
\end{tikzpicture}
\end{center}

\subsubsection{Bus parallèle}

\begin{defbox}{Définition : Bus parallèle}
On parle de communication et de \textbf{bus parallèle} quand les bits circulent \textbf{ensemble} sur des supports séparés.

\vspace{0.3cm}
\textbf{Exemples :} bus mémoire, Centronic, PATA, SCSI, PCI
\end{defbox}

\begin{center}
\textbf{Représentation d'un bus parallèle}

\vspace{0.5cm}
\begin{tikzpicture}
\foreach \y/\bit in {0/1, 1/0, 2/1, 3/1, 4/0, 5/0, 6/1, 7/0} {
    \draw[thick, ->] (0,\y*0.5) -- (5,\y*0.5);
    \node[left] at (-0.3,\y*0.5) {\small \bit};
    \node[right] at (5.3,\y*0.5) {\small \bit};
}
\node[above] at (2.5,4) {10110010};
\end{tikzpicture}
\end{center}

\subsubsection{Comparaison}

\begin{table}[h]
\centering
\begin{tabular}{|l|c|c|}
\hline
\textbf{Critère} & \textbf{Bus Série} & \textbf{Bus Parallèle} \\
\hline
Nombre de fils & Peu & Nombreux \\
\hline
Coût & Moins cher & Plus cher \\
\hline
Vitesse (théorique) & Plus lent & Plus rapide \\
\hline
Exemples & USB, SATA & DATA, PCI, IDE \\
\hline
\end{tabular}
\caption{Comparaison bus série vs bus parallèle}
\end{table}

\begin{remarkbox}{Remarque importante}
Paradoxalement, avec les technologies modernes, les bus série sont devenus très rapides (USB 3.0, SATA III) et ont remplacé progressivement les bus parallèles qui souffraient de problèmes de synchronisation à haute fréquence.
\end{remarkbox}

\subsection{Aiguillage de l'information}

\begin{defbox}{Principe de l'aiguillage}
Il n'y a pas un bus par couple de circuits à connecter, mais souvent plusieurs circuits se partagent un bus pour échanger entre eux. 

L'information qui circule sur un bus doit donc être :
\begin{itemize}
    \item \textbf{Sélectionnée} (source)
    \item \textbf{Dirigée} (destination)
\end{itemize}

Il faut donc un moyen de faire passer ou d'empêcher le passage d'une information sur un fil.
\end{defbox}

\subsubsection{Autorisation de passage par porte ET}

On peut empêcher le passage d'une information sur un fil en utilisant une porte ET, en disant que l'absence d'information correspond à la valeur 0.

\begin{center}
\textbf{Notations :}
\begin{itemize}
    \item $E$ : information en entrée
    \item $S$ : information en sortie
    \item $C$ : commande de passage
\end{itemize}
\end{center}

\begin{table}[h]
\centering
\begin{tabular}{|c|c|c|}
\hline
$C$ & $E$ & $S = C \cdot E$ \\
\hline
0 & 0 & 0 \\
0 & 1 & 0 \\
1 & 0 & 0 \\
1 & 1 & 1 \\
\hline
\end{tabular}
\end{table}

\begin{propbox}{Interprétation}
\begin{center}
\begin{tabular}{|c|c|}
\hline
$C$ & $S$ \\
\hline
0 & 0 \\
1 & $E$ \\
\hline
\end{tabular}
\end{center}

\begin{itemize}
    \item Si $C = 0$ alors la sortie $S$ vaut 0 quelle que soit la valeur de $E$
    \item Si $C = 1$ alors la sortie $S$ vaut $E$
\end{itemize}
\end{propbox}

\begin{center}
\begin{circuitikz}
\draw (0,2) node[and port] (and1) {};
\draw (and1.in 1) -- ++(-1,0) node[left] {$C$};
\draw (and1.in 2) -- ++(-1,0) node[left] {$E$};
\draw (and1.out) -- ++(1,0) node[right] {$S$};
\end{circuitikz}
\end{center}

\subsubsection{Autorisation de passage par porte OU}

On peut empêcher le passage d'une information sur un fil en utilisant une porte OU, en disant que l'absence d'information correspond à la valeur 1.

\begin{table}[h]
\centering
\begin{tabular}{|c|c|c|}
\hline
$C$ & $E$ & $S = C + E$ \\
\hline
0 & 0 & 0 \\
0 & 1 & 1 \\
1 & 0 & 1 \\
1 & 1 & 1 \\
\hline
\end{tabular}
\end{table}

\begin{propbox}{Interprétation}
\begin{center}
\begin{tabular}{|c|c|}
\hline
$C$ & $S$ \\
\hline
0 & $E$ \\
1 & 1 \\
\hline
\end{tabular}
\end{center}

\begin{itemize}
    \item Si $C = 1$ alors la sortie $S$ vaut 1 quelle que soit la valeur de $E$
    \item Si $C = 0$ alors la sortie $S$ vaut $E$
\end{itemize}
\end{propbox}

\newpage
\section{Le Multiplexeur}

\subsection{Description du circuit}

\begin{defbox}{Définition : Multiplexeur (MUX)}
Un \textbf{multiplexeur} est un circuit permettant de faire passer sur un bus l'information en provenance d'un autre bus choisi parmi plusieurs.

Pour sélectionner le bus source, il faut lui fournir son numéro ou, en terme plus approprié, son \textbf{adresse}.
\end{defbox}

\subsubsection{Composants d'entrée}

Le multiplexeur comprend en entrée :
\begin{itemize}
    \item \textbf{$N$ bus de données} identiques, comportant le même nombre de ``fils'', nommés $B_0, B_1, \ldots, B_{N-1}$
    
    Les fils du bus $B_i$ seront eux aussi numérotés : $B_{i,0}, B_{i,1}, \ldots, B_{i,M-1}$, où $M$ est le nombre de fils de chacun des bus.
    
    \item Un \textbf{bus de sélection} ou \textbf{bus d'adresse} de $n$ fils tel que $N = 2^n$, nommé $A$
    
    \item Une \textbf{commande d'activation} $CS$ (Chip Select)
\end{itemize}

\subsubsection{Composants de sortie}

En sortie, un bus de données identique aux bus de données d'entrée, nommé $S$, comprenant $M$ fils eux-mêmes nommés $S_0, S_1, \ldots, S_{M-1}$.

\subsection{Fonctionnement}

\begin{propbox}{Règle de fonctionnement}
\begin{itemize}
    \item Si $CS = 0$ alors $S = 0$
    \item Si $CS = 1$ et $A = (i)_2$ alors $S = B_i$
\end{itemize}
\end{propbox}

\begin{center}
\begin{tikzpicture}[scale=0.9]
\draw[thick] (0,0) rectangle (3,5);
\node at (1.5,2.5) {\Large MUX};

% Entrées
\draw[thick, <-] (0,4.5) -- (-1,4.5) node[left] {$B_0$};
\draw[thick, <-] (0,3.5) -- (-1,3.5) node[left] {$B_1$};
\node at (-0.5,2.5) {$\vdots$};
\draw[thick, <-] (0,0.5) -- (-1,0.5) node[left] {$B_{N-1}$};

% Sortie
\draw[thick, ->] (3,2.5) -- (4,2.5) node[right] {$S$};

% Commandes
\draw[thick, <-] (1.5,5) -- (1.5,5.8) node[above] {$CS$};
\draw[thick, <-] (0.5,5) -- (0.5,6.3) node[above] {$A$};
\end{tikzpicture}
\end{center}

\subsection{Réalisation du circuit}

Il faut envoyer vers la sortie $S$ l'information portée par le bus dont le numéro $i$ est mis sur le bus d'adresse. Cette adresse va être décodée par un décodeur.

\begin{examplebox}{Exemple : Multiplexeur 4 vers 1 (8 fils)}
Pour un multiplexeur de 4 bus de 8 fils de large :
\begin{itemize}
    \item Pour choisir parmi 4 bus il faut 4 valeurs
    \item Donc 2 bits d'adresse
    \item Donc un décodeur $2 \times 4$
    \item Les 4 bus d'entrée et le bus de sortie ont 8 fils chacun
\end{itemize}

Le schéma utilise des portes ET pour chaque fil, activées par les sorties du décodeur qui analyse l'adresse, puis une porte OU pour combiner les résultats.
\end{examplebox}

\subsection{Applications des multiplexeurs}

\subsubsection{Génération de fonction logique}

\begin{propbox}{Principe}
Une fonction logique définit une valeur binaire pour toutes les combinaisons d'un ensemble de variables binaires. 

Si on présente ces valeurs en entrée d'un multiplexeur à un bit de donnée, et qu'on présente les variables sur le bus d'adresse, on retrouve en sortie la valeur de la fonction.
\end{propbox}

\begin{examplebox}{Exemple : Fonction à 3 variables}
Soit la fonction :
\[f_2(a,b,c) = \overline{a} \cdot \overline{b} \cdot \overline{c} + \overline{a} \cdot b \cdot \overline{c} + a \cdot \overline{b} \cdot \overline{c} + a \cdot \overline{b} \cdot c + a \cdot b \cdot c\]

\begin{center}
\begin{tabular}{|c|c|c|c|}
\hline
$a$ & $b$ & $c$ & $f_2$ \\
\hline
0 & 0 & 0 & 1 \\
0 & 0 & 1 & 0 \\
0 & 1 & 0 & 1 \\
0 & 1 & 1 & 0 \\
1 & 0 & 0 & 1 \\
1 & 0 & 1 & 1 \\
1 & 1 & 0 & 0 \\
1 & 1 & 1 & 1 \\
\hline
\end{tabular}
\end{center}

En utilisant un multiplexeur à 3 bits d'adresse avec les constantes 1,0,1,0,1,1,0,1 en entrée, on obtient un circuit équivalent à la fonction.

\textbf{Intérêt :} La fonction est paramétrable par les constantes présentées en entrée du multiplexeur.
\end{examplebox}

\subsubsection{Conversion parallèle/série}

\begin{propbox}{Principe}
La coexistence dans un ordinateur de plusieurs types de bus rend nécessaire diverses conversions. 

Pour sérialiser une donnée fournie par un bus parallèle de $N$ bits, il faut :
\begin{itemize}
    \item Un multiplexeur à $N$ bus d'entrée de 1 bit de donnée
    \item Envoyer la donnée sur les bus d'entrée, 1 bit par bus
    \item Envoyer successivement sur le bus d'adresse les valeurs croissantes à partir de 0
\end{itemize}

Les bits de la donnée se trouveront successivement sur la sortie.
\end{propbox}

\newpage
\section{Le Démultiplexeur}

\subsection{Description}

\begin{defbox}{Définition : Démultiplexeur (DEMUX)}
Un \textbf{démultiplexeur} est le circuit réalisant l'opération inverse du multiplexeur.

Il prend une information sur un bus d'entrée et la dirige vers l'un des $N$ bus de sortie selon l'adresse fournie.
\end{defbox}

\subsection{Composants}

Ce circuit comporte :
\begin{itemize}
    \item Un \textbf{bus en entrée} de $M$ fils parallèles, nommé $B$
    \item $N$ \textbf{bus de sortie} identiques au bus d'entrée, nommés $S_0, S_1, \ldots, S_{N-1}$
    \item Un \textbf{bus d'adresse} de $n$ bits tel que $N = 2^n$, nommé $A$
    \item Une \textbf{commande d'activation} $CS$
\end{itemize}

\begin{center}
\begin{tikzpicture}[scale=0.9]
\draw[thick] (0,0) rectangle (3,5);
\node at (1.5,2.5) {\Large DEMUX};

% Entrée
\draw[thick, <-] (0,2.5) -- (-1,2.5) node[left] {$B$};

% Sorties
\draw[thick, ->] (3,4.5) -- (4,4.5) node[right] {$S_0$};
\draw[thick, ->] (3,3.5) -- (4,3.5) node[right] {$S_1$};
\node at (3.5,2.5) {$\vdots$};
\draw[thick, ->] (3,0.5) -- (4,0.5) node[right] {$S_{N-1}$};

% Commandes
\draw[thick, <-] (1.5,5) -- (1.5,5.8) node[above] {$CS$};
\draw[thick, <-] (0.5,5) -- (0.5,6.3) node[above] {$A$};
\end{tikzpicture}
\end{center}

\subsection{Fonctionnement}

\begin{propbox}{Règle de fonctionnement}
\begin{itemize}
    \item Si $CS = 0$ alors toutes les sorties valent 0 : $S_i = 0, \forall i$
    \item Si $CS = 1$ et $A = (i)_2$ alors $S_i = B$ et $S_j = 0, \forall j \neq i$
\end{itemize}
\end{propbox}

\subsection{Réalisation}

Le principe est similaire au multiplexeur mais inversé :
\begin{itemize}
    \item Un décodeur analyse l'adresse $A$
    \item Chaque sortie du décodeur active un ensemble de portes ET
    \item Seul le bus de sortie correspondant à l'adresse reçoit l'information
\end{itemize}

\subsection{Application : Conversion série/parallèle}

Le démultiplexeur permet de convertir une information série en parallèle :
\begin{itemize}
    \item L'information série arrive bit par bit sur l'entrée $B$
    \item L'adresse change séquentiellement
    \item Chaque bit est dirigé vers sa position dans le bus parallèle de sortie
\end{itemize}

\newpage
\section{L'Unité Arithmétique et Logique (UAL/ALU)}

\subsection{Présentation générale}

\begin{defbox}{Définition : ALU}
L'\textbf{Unité Arithmétique et Logique} (UAL en français, ALU en anglais pour \textit{Arithmetic Logic Unit}) est le composant du processeur qui effectue les opérations arithmétiques et logiques.

Elle prend en entrée des opérandes et produit un résultat ainsi que des \textbf{flags} (drapeaux) indiquant des informations sur le résultat de l'opération.
\end{defbox}

\subsection{Architecture de l'ALU}

\begin{center}
\begin{tikzpicture}[scale=1]
\draw[thick] (0,0) rectangle (5,4);
\node at (2.5,2) {\Large \textbf{ALU}};

% Entrées
\draw[thick, <-] (0,3) -- (-2,3) node[left] {Opérande A};
\draw[thick, <-] (0,1) -- (-2,1) node[left] {Opérande B};
\draw[thick, <-] (2.5,0) -- (2.5,-1) node[below] {Opération};

% Sorties
\draw[thick, ->] (5,2.5) -- (7,2.5) node[right] {Résultat};
\draw[thick, ->] (5,1.5) -- (7,1.5) node[right] {Flags (ZF, SF, CF, OF, PF)};
\end{tikzpicture}
\end{center}

\subsection{Les flags (drapeaux)}

Les flags sont des bits qui donnent des informations sur le résultat de l'opération effectuée par l'ALU. Ils sont essentiels pour les instructions conditionnelles et la gestion des erreurs.

\subsubsection{Zero Flag (ZF)}

\begin{defbox}{Définition : Zero Flag (ZF)}
Le \textbf{Zero Flag} indique si le résultat de l'opération est nul.

\[ZF = \begin{cases}
1 & \text{si le résultat} = 0 \\
0 & \text{sinon}
\end{cases}\]

Formule logique : $ZF = \overline{a_0 + a_1 + a_2 + \ldots + a_n}$

où $a_i$ sont les bits du résultat (on fait un OU de tous les bits, puis on nie).
\end{defbox}

\begin{examplebox}{Exemple}
\begin{itemize}
    \item $5 - 5 = 0$ $\Rightarrow$ ZF = 1
    \item $7 - 3 = 4 = 0100_2$ $\Rightarrow$ ZF = 0
\end{itemize}
\end{examplebox}

\subsubsection{Sign Flag (SF)}

\begin{defbox}{Définition : Sign Flag (SF)}
Le \textbf{Sign Flag} indique le signe du résultat en représentation complément à deux.

\[SF = \text{bit de poids fort du résultat}\]

\begin{itemize}
    \item SF = 0 $\Rightarrow$ résultat positif
    \item SF = 1 $\Rightarrow$ résultat négatif
\end{itemize}
\end{defbox}

\begin{examplebox}{Exemple (sur 8 bits)}
\begin{itemize}
    \item $00101100_2 = 44_{10}$ $\Rightarrow$ SF = 0 (positif)
    \item $11010011_2 = -45_{10}$ (C2) $\Rightarrow$ SF = 1 (négatif)
\end{itemize}
\end{examplebox}

\subsubsection{Carry Flag (CF)}

\begin{defbox}{Définition : Carry Flag (CF)}
Le \textbf{Carry Flag} indique s'il y a eu une \textbf{retenue} (carry) ou un \textbf{emprunt} (borrow) lors de l'opération sur le bit de poids fort.

\begin{itemize}
    \item Pour une addition : CF = 1 s'il y a retenue sortante
    \item Pour une soustraction : CF = 1 s'il y a emprunt
\end{itemize}

Le CF est utilisé pour l'arithmétique \textbf{non signée}.
\end{defbox}

\begin{examplebox}{Exemple : Addition (sur 4 bits)}
\begin{align*}
  &\phantom{+} 1011_2 \quad (11_{10}) \\
+ &\phantom{+} 0111_2 \quad (7_{10}) \\
\hline
&\textcolor{red}{1}\, 0010_2 \quad (2_{10} \text{ sur 4 bits})
\end{align*}

Retenue sortante $\Rightarrow$ CF = 1

Le résultat devrait être 18 mais sur 4 bits on obtient 2 (débordement non signé).
\end{examplebox}

\subsubsection{Overflow Flag (OF)}

\begin{defbox}{Définition : Overflow Flag (OF)}
L'\textbf{Overflow Flag} indique un \textbf{débordement en arithmétique signée} (complément à deux).

Il signale que le résultat est trop grand (ou trop petit) pour être représenté correctement dans le nombre de bits disponible.

\textbf{Règle de détection :}
\[OF = C_{in}(\text{MSB}) \oplus C_{out}(\text{MSB})\]

où $C_{in}$ est la retenue entrante sur le bit de poids fort et $C_{out}$ la retenue sortante.
\end{defbox}

\begin{examplebox}{Exemple : Addition signée (sur 4 bits)}
\begin{align*}
  &\phantom{+} 0110_2 \quad (+6_{10}) \\
+ &\phantom{+} 0101_2 \quad (+5_{10}) \\
\hline
&\phantom{+} 1011_2 \quad (-5_{10} \text{ en C2, mais devrait être } +11)
\end{align*}

Deux nombres positifs donnent un résultat négatif $\Rightarrow$ OF = 1

Le résultat est incorrect en arithmétique signée.
\end{examplebox}

\begin{remarkbox}{Différence CF vs OF}
\begin{itemize}
    \item \textbf{CF} : pour l'arithmétique \textbf{non signée} (entiers naturels)
    \item \textbf{OF} : pour l'arithmétique \textbf{signée} (entiers relatifs en C2)
\end{itemize}
\end{remarkbox}

\subsubsection{Parity Flag (PF)}

\begin{defbox}{Définition : Parity Flag (PF)}
Le \textbf{Parity Flag} indique la \textbf{parité} du nombre de bits à 1 dans le résultat.

\[PF = \bigoplus_{i=0}^{n-1} a_i\]

(XOR de tous les bits du résultat)

\begin{itemize}
    \item PF = 0 $\Rightarrow$ nombre pair de bits à 1 (parité paire)
    \item PF = 1 $\Rightarrow$ nombre impair de bits à 1 (parité impaire)
\end{itemize}
\end{defbox}

\begin{examplebox}{Exemple}
\begin{itemize}
    \item $10110010_2$ : 4 bits à 1 $\Rightarrow$ PF = 0 (pair)
    \item $10110110_2$ : 5 bits à 1 $\Rightarrow$ PF = 1 (impair)
\end{itemize}

Calcul pour $1011_2$ :
\[PF = 1 \oplus 0 \oplus 1 \oplus 1 = 1\]
\end{examplebox}

\subsection{Tableau récapitulatif des flags}

\begin{table}[h]
\centering
\begin{tabular}{|l|l|p{7cm}|}
\hline
\textbf{Flag} & \textbf{Nom} & \textbf{Signification} \\
\hline
ZF & Zero Flag & Résultat = 0 \\
\hline
SF & Sign Flag & Signe du résultat (bit de poids fort) \\
\hline
CF & Carry Flag & Retenue/emprunt (arithmétique non signée) \\
\hline
OF & Overflow Flag & Débordement (arithmétique signée) \\
\hline
PF & Parity Flag & Parité du résultat (XOR des bits) \\
\hline
\end{tabular}
\caption{Les principaux flags de l'ALU}
\end{table}

\subsection{Utilisation des flags}

Les flags sont utilisés par les instructions conditionnelles pour prendre des décisions :

\begin{examplebox}{Exemples d'utilisation}
\begin{itemize}
    \item \textbf{JZ} (Jump if Zero) : sauter si ZF = 1
    \item \textbf{JS} (Jump if Sign) : sauter si SF = 1 (négatif)
    \item \textbf{JC} (Jump if Carry) : sauter si CF = 1
    \item \textbf{JO} (Jump if Overflow) : sauter si OF = 1
    \item \textbf{JP} (Jump if Parity) : sauter si PF = 1 (parité impaire)
    \item \textbf{JNZ} (Jump if Not Zero) : sauter si ZF = 0
\end{itemize}
\end{examplebox}

\begin{remarkbox}{Application pratique}
Les flags permettent au processeur de :
\begin{itemize}
    \item Détecter les débordements arithmétiques
    \item Implémenter des boucles (comparaison avec ZF)
    \item Gérer les nombres signés et non signés
    \item Détecter des erreurs de transmission (PF)
    \item Prendre des décisions conditionnelles dans le code machine
\end{itemize}
\end{remarkbox}

\newpage
\section{Conclusion et points clés}

\subsection{Optimisation des circuits}

L'utilisation des états indéterminés permet de :
\begin{itemize}
    \item Simplifier les fonctions logiques
    \item Économiser des bits
    \item Réduire la complexité des circuits
    \item Optimiser l'espace physique sur les circuits intégrés
\end{itemize}

\subsection{Choix du système de codage}

Le choix du système de codage dépend de l'application :
\begin{itemize}
    \item \textbf{BCD/Aiken} : interfaçage avec des afficheurs, calculs décimaux
    \item \textbf{Code de GRAY} : encodeurs rotatifs, compteurs fiables
    \item \textbf{Excédant 3} : détection d'erreurs, arithmétique
    \item \textbf{ASCII} : représentation de caractères et texte
\end{itemize}

\subsection{Architecture et circuits}

Les circuits de transcodage (codeur, décodeur, transcodeur) sont essentiels pour :
\begin{itemize}
    \item Interfacer différents systèmes
    \item Optimiser la représentation de l'information
    \item Faciliter les calculs et opérations
    \item Permettre la détection d'erreurs
\end{itemize}

\subsection{Circulation de l'information}

Les multiplexeurs et démultiplexeurs permettent :
\begin{itemize}
    \item De partager efficacement les bus entre plusieurs circuits
    \item De convertir entre formats série et parallèle
    \item De générer des fonctions logiques configurables
    \item D'optimiser le routage des données dans le processeur
\end{itemize}

\subsection{L'ALU et les flags}

L'ALU est le cœur du processeur pour les opérations arithmétiques et logiques :
\begin{itemize}
    \item Les flags fournissent des informations cruciales sur les résultats
    \item ZF, SF, CF, OF et PF permettent la gestion des conditions
    \item Distinction importante entre arithmétique signée (OF) et non signée (CF)
    \item Les flags sont la base des instructions conditionnelles
\end{itemize}

\vspace{2cm}

\begin{center}
\rule{10cm}{0.5pt}

\vspace{0.5cm}

\textit{Ce document présente les fondements de l'algèbre de Boole, \\
des fonctions logiques, de la représentation de l'information, \\
de la circulation de l'information dans les circuits \\
et de l'architecture des systèmes numériques.}

\vspace{0.5cm}

\textbf{Arthur LOUVET - Licence Informatique}

\textbf{Semestres 1 \& 2 (2024-2025)}
\end{center}

\end{document}